\documentclass[10pt]{beamer}
\usetheme{Madrid}
\usepackage[T1]{fontenc}
\usepackage[utf8]{inputenc}
\usepackage[brazil]{babel}
\usepackage{lmodern}
\usepackage{hyperref}
\usepackage{listings}
\usepackage{xcolor}
\setbeamertemplate{navigation symbols}{}

\lstset{
  language=C++,
  basicstyle=\ttfamily\small,
  keywordstyle=\color{blue},
  commentstyle=\color{gray},
  stringstyle=\color{red},
  breaklines=true,
  showstringspaces=false,
}

\title{Bibliotecas Comuns}
\subtitle{Curso de Microcontroladores}
\author{}
\date{}

\begin{document}

\begin{frame}
  \titlepage
\end{frame}

\begin{frame}{O que é uma biblioteca?}
  \begin{itemize}
    \item Código pronto para usar de novo e de novo.
    \item Implementa funções complexas de forma simples.
    \item Economiza tempo e reduz erros.
    \item Arduino vem com muitas bibliotecas prontas.
  \end{itemize}
\end{frame}

\begin{frame}[fragile]{Incluindo uma biblioteca}
\begin{lstlisting}
#include <Servo.h>
#include <Wire.h>
#include <LiquidCrystal.h>
\end{lstlisting}

  \begin{itemize}
    \item `#include` traz a biblioteca para o seu sketch.
    \item Use `<...>` para bibliotecas do sistema.
    \item Use `\"...\"` para bibliotecas locais.
  \end{itemize}
\end{frame}

\begin{frame}[fragile]{Biblioteca Servo - exemplo}
\begin{lstlisting}
#include <Servo.h>

Servo meuServo;
const int PINO_SERVO = 9;

void setup() {
  meuServo.attach(PINO_SERVO);
}

void loop() {
  meuServo.write(0);   // 0 graus
  delay(1000);
  meuServo.write(90);  // 90 graus
  delay(1000);
  meuServo.write(180); // 180 graus
  delay(1000);
}
\end{lstlisting}
\end{frame}

\begin{frame}{Biblioteca Servo - uso na prática}
  \begin{itemize}
    \item `attach()` conecta o servo a um pino com PWM.
    \item `write()` define o ângulo em graus (0-180).
    \item Servo controlado por pulsos, a biblioteca gerencia tudo.
    \item Essencial para projetos com braços robóticos.
  \end{itemize}
\end{frame}

\begin{frame}{Outras bibliotecas importantes}
  \begin{itemize}
    \item `Wire`: I2C para sensores IMU, acelerômetro, etc.
    \item `SPI`: comunicação rápida com cartão SD, display.
    \item `LiquidCrystal`: controla displays LCD de 16x2.
    \item `Adafruit_*`: bibliotecas específicas de sensores.
  \end{itemize}
\end{frame}

\begin{frame}{Como encontrar e instalar}
  \begin{itemize}
    \item Arduino IDE tem Sketch > Include Library > Manage Libraries.
    \item Procure por nome da biblioteca ou do sensor.
    \item Exemplos vêm junto com a instalação.
    \item Leia o arquivo README para instruções.
  \end{itemize}
\end{frame}

\begin{frame}{Estruturando código com bibliotecas}
  \begin{itemize}
    \item Separe lógica de sensores em funções simples.
    \item Use classes para agrupar comportamento (objeto-orientado).
    \item Mantenha sketch principal limpo e organizado.
    \item Comente o que cada função faz.
  \end{itemize}
\end{frame}

\end{document}
